%% Mẫu chung cho mọi approach. File approach_template.tex không được biên dịch.
%% Thêm approach mới: sao chép mẫu này và thêm \input trong notes.tex.
%% Chỉ giữ đúng thứ tự các mục. Không có dữ liệu thì ghi Not Evaluated hoặc Not Supported.
\chapter{Approach X: [TBD]}

\begin{table}[H]
  \centering
  \caption{Thông tin approach X}
  \begin{tabularx}{\linewidth}{l X}
    \toprule
    \textbf{Field} & \textbf{Giá trị} \\
    \midrule
    Approach name & TBD \\
    Paper / Repository & TBD \\
    Owner & TBD \\
    Track & A / B / Both \\
    Camera model & TBD \\
    LiDAR support & Yes / No / Optional \\
    Required inputs & TBD \\
    Output representation & TBD \\
    GPU requirements & TBD \\
    Reproduction status & Not Started \\
    CARLA feasibility & Not Evaluated \\
    \bottomrule
  \end{tabularx}
\end{table}

\section{Overview \& References}

Hướng làm: ghi tên paper, repository, tác giả, mục tiêu phương pháp, phiên bản và trạng thái mã nguồn.

\section{Methodology \& Technical Background}

Hướng làm: trình bày ý tưởng cốt lõi, thuật toán chính, assumptions và khác biệt với baseline.

\section{Original Pipeline Architecture}

Hướng làm: vẽ pipeline nguyên bản, ghi từng stage, input/output và dependencies.
Mô tả đúng pipeline thực tế của approach, không ép theo pipeline chung.

\begin{table}[H]
  \centering
  \caption{Ánh xạ pipeline của approach X sang module của framework}
  \begin{tabularx}{\linewidth}{X X l}
    \toprule
    \textbf{Stage gốc} & \textbf{Module của framework} & \textbf{Thay thế được} \\
    \midrule
    TBD & TBD & Yes / No \\
    \bottomrule
  \end{tabularx}
\end{table}

\section{Requirements \& Environment Setup}

Hướng làm: ghi camera models, LiDAR, calibration, GPU, VRAM, software và cách cài đặt.

\section{Implementation \& Reproduction}

Hướng làm: ghi cách chạy thật, preprocessing, config, các thay đổi trong code và các bước reproduce.

\section{Experiments \& Configuration}

Hướng làm: ghi dataset, sequences, train/test split, parameters và experiment IDs.

\section{Results \& Benchmark}

Hướng làm: đưa ảnh render, reconstructed scene, số đo và logs.

\begin{table}[H]
  \centering
  \caption{Kết quả benchmark của approach X}
  \begin{tabularx}{\linewidth}{X l l}
    \toprule
    \textbf{Metric} & \textbf{Experiment ID} & \textbf{Kết quả} \\
    \midrule
    PSNR & TBD & Not Evaluated \\
    SSIM & TBD & Not Evaluated \\
    LPIPS & TBD & Not Evaluated \\
    Geometry & TBD & Not Evaluated \\
    Runtime & TBD & Not Evaluated \\
    VRAM & TBD & Not Evaluated \\
    \bottomrule
  \end{tabularx}
\end{table}

\section{Analysis \& Failure Cases}

Hướng làm: phân tích kết quả tốt và xấu, nguyên nhân, giới hạn và trade-offs.

\section{Framework \& CARLA Compatibility}

Hướng làm: ghi module tái sử dụng được, artifacts, adapter cần thiết và CARLA feasibility.

\section{Conclusions \& Reproducibility}

Hướng làm: nêu kết luận, source commit, config, commands và artifacts bàn giao.
